\SourceLeaf{304}
\Needspace{6\baselineskip}\begin{center}\textsc{Gauss} to \textsc{Schumacher}. Göttingen, 15 May 1843.\end{center}

I add a few words to my last letter to you, my dearest friend,
to do justice to a book to which, partly because of a circumstance
you brought about twenty years ago (quite innocently, of course),
I have so far paid less attention than it deserves.

You may still remember that in your house at the time you presented me with a rather
elegant geometrical problem posed by \textsc{Möbius}, for which I
almost immediately wrote down a solution that left nothing to be desired. As a result
of this incident, I took up \textsc{Möbius}'s book, barycentrischer Calcul,
in whose preface that problem was also mentioned, with a kind of prejudice.
I doubted whether it was worth the trouble to master a rather ingeniously
devised method of calculation if it could achieve nothing that could not
be achieved just as easily without it. I had therefore laid the
book aside for the time being --- especially since I received it when other things
fully occupied me --- and, as has happened to me with many things, had expected little
of it and later forgotten it altogether.

In the last few days (yesterday, to be exact) I happened to come across the book,
and I soon found with great pleasure that it contains the quintessence
of the theory of conics in a nutshell, and that his
barycentric calculus leads by the easiest route to the solution of all the relevant
problems. While leafing through it, I found much else of interest as well,
for example the beautiful theorem discovered by \textsc{Mac Laurin} and \textsc{Brackenridge},
which I had either never known or forgotten again. By means of it
the problem of drawing a conic through 5 given points can be solved so
elegantly.
